\chapter{Why This Stack}
\label{ch:why}

This book documents the design, theory, and hard lessons of the Vanguard rover arm ---
manipulation autonomy for the BITS Pilani Hyderabad Mars Rover team, built to win IRC~2027 and
grow into a research platform. Every chapter is written in the same session the work happened,
so the reasoning here is what we actually thought, not a cleaned-up retrospective.

\section{The mission, worked backwards}

The arm exists to score points in four competitions (IRC, ARC, ERC, URC) whose manipulation
tasks --- studied from the primary rulebooks --- collapse into six reusable skills: grasp \&
carry, panel operations, connector insertion, oriented placement, sampling support, and
read-\&-report. IRC's IDMO panel (buttons, switches, knobs, a 3-pin plug), ARC's propellant
hose, ERC's RJ-45, and URC's board connectors are one skill family wearing four costumes. The
central architectural bet follows: \emph{build skills, not missions}. A new rulebook becomes a
re-orchestration, not a rebuild.

Two strategy decisions bound everything else:
\begin{enumerate}
  \item \textbf{Reliable assisted teleoperation wins IRC 2027; genuine autonomy is the ERC 2027
  goal.} Top teams win on reliability and operator skill. Autonomy is invested exactly where the
  rulebook prices it (IRC's RADO delivery: autonomous = 100\% of points, teleoperated = 50\%).
  \item \textbf{Simulation first.} The physical arm arrives months after work begins. The stack
  is therefore built against an Isaac Sim digital twin exposing \emph{identical ROS~2
  interfaces} to the future hardware --- so hardware arrival is an integration event, not a
  starting gun. This same twin later becomes the benchmark environment of our first paper.
\end{enumerate}

\section{One mental model}
\label{sec:mental-model}

A manipulator is a chain of rigid-body transforms in $SE(3)$, and every tool in the stack is
that one mathematical object viewed differently. The URDF \emph{declares} the chain. TF2
\emph{broadcasts} it. Forward kinematics \emph{composes} it:
\[
  T_{\text{ee}}(\theta) \;=\; e^{[\mathcal{S}_1]\theta_1} e^{[\mathcal{S}_2]\theta_2} \cdots
  e^{[\mathcal{S}_n]\theta_n}\, M ,
\]
the product of exponentials over the joint screw axes $\mathcal{S}_i$ acting on the home pose
$M$. Inverse kinematics \emph{inverts} it --- nonlinear, multi-solution, and ill-conditioned
near singularities. The Jacobian \emph{differentiates} it, $\mathcal{V} = J(\theta)\dot\theta$,
and is why teleop jogging, visual servoing, and singularity handling are the same problem.
MoveIt~2 \emph{plans} through the chain; \texttt{ros2\_control} \emph{executes} along it; Isaac
Sim \emph{simulates} it from the same URDF. Ten tools, one idea.

Two consequences we will meet repeatedly:
\begin{itemize}
  \item \textbf{Six degrees of freedom is the floor, not a luxury.} Connector insertion
  requires commanding an arbitrary approach \emph{orientation} at an arbitrary position ---
   six constraints. A 5-DoF arm can reach the socket but cannot, in general, present the plug
  square to it.
  \item \textbf{Singularity is a property of the map, not the motor.} At a straight elbow,
  $J$ loses rank; some Cartesian velocities become unreachable and naive IK commands explode.
  We damp (damped least squares) rather than pretend precision the hardware doesn't have.
\end{itemize}

\section{Infrastructure decisions (recorded before the first line of code)}
\label{sec:infra}

\begin{description}
  \item[Platform.] Ubuntu 22.04 + ROS~2 Humble for the IRC season; migration to Jazzy/24.04 is
  \emph{scheduled} (post-IRC window, Humble EOL May 2027) rather than improvised.
  \item[Containers.] Development happens in a devcontainer (\texttt{ros:humble-ros-base-jammy});
  the rover will run a lean L4T-based runtime image on the Jetson. Deployment is
  \texttt{git pull \&\& docker compose up}, not dependency archaeology at 7\,am in a field.
  \item[DDS.] CycloneDDS as the team RMW, \texttt{--network=host --ipc=host} on every container,
  one \texttt{ROS\_DOMAIN\_ID}. This dissolves the classic ``container can't discover the
  robot'' failure class before anyone hits it.
  \item[TensorRT engines are built on the device that runs them.] Engines are not portable
  artifacts; on-rover engine generation is a scripted first-boot step (a lesson inherited
  directly from the APEX project).
  \item[Bring-up is native; runtime is containerized.] When a motor doesn't enumerate you want
  \texttt{dmesg} with zero layers in between. Containers earn their keep once drivers are
  stable.
  \item[Compute.] Development and Isaac scenes on a 12\,GB RTX laptop; heavy training
  (GR00T-class fine-tunes, RL) as containerized batch jobs on the university HPC. The nightly
  simulation regression runs on the docked laptop.
\end{description}

\section{Mistakes log, opened early}
\label{sec:mistakes-ch1}

This section exists in every chapter. Failures are recorded with root causes because they are
the highest-density teaching material we produce.

\begin{lesson}
\textbf{The paste-boundary error (day one).} The very first typed file of the project ---
\texttt{.devcontainer/Dockerfile.dev} --- ended up containing the JSON of
\texttt{devcontainer.json} pasted below the Dockerfile instructions; the build would have died
on the first JSON line. Found in review before first build. Root cause: transcribing two
reference files in one editing pass without a file-boundary check. Rule adopted: after typing
from reference, \emph{verify each file builds/parses on its own} before moving to the next.
The same failure family (paste-indent) cost a debugging session on APEX in June 2026.
\end{lesson}
